\documentclass[conference]{IEEEtran}
\IEEEoverridecommandlockouts
% The preceding line is only needed to identify funding in the first footnote. If that is unneeded, please comment it out.
\usepackage{cite}
\usepackage{amsmath,amssymb,amsfonts}
\usepackage{algorithmic}
\usepackage{graphicx}
\usepackage{textcomp}
\usepackage{xcolor}
\usepackage{booktabs}
\usepackage{tabularx}
\usepackage{multirow}
\usepackage{enumitem}
\usepackage{url}
\usepackage{hyperref}
\usepackage{microtype}
\usepackage{algorithm}

\def\BibTeX{{\rm B\kern-.05em{\sc i\kern-.025em b}\kern-.08em
    T\kern-.1667em\lower.7ex\hbox{E}\kern-.125emX}}

\newtheorem{definition}{Definition}
\newtheorem{theorem}{Theorem}

\begin{document}

\title{Officer's Arena: An Adaptive Multidimensional Psychometric Engine and Multimodal Retrieval-Augmented System for High-Stakes Competitive Examinations\\
}

\author{\IEEEauthorblockN{1\textsuperscript{st} Pranay Sharma}
\IEEEauthorblockA{\textit{Scope of Computer Science and Engineering(SCOPE)} \\
\textit{Vellore Institute of Technology}\\
Chennai, India \\
pranay.sharma2022@vitstudent.ac.in}
\and
\IEEEauthorblockN{2\textsuperscript{nd} Rathna R}
\IEEEauthorblockA{\textit{Scope of Computer Science and Engineering(SCOPE)} \\
\textit{Vellore Institute of Technology}\\
Chennai, India \\
rathna.r@vit.ac.in}
}

\maketitle

\begin{abstract}
High-stakes standardized competitive examinations---exemplified by the Union Public Service Commission Civil Services Examination (UPSC CSE) and Combined Defence Services (CDS) examination---demand extensive multi-disciplinary domain synthesis, high-order cognitive reasoning, and rigorous long-term knowledge retention across thousands of syllabus sub-facets. Traditional computer-based preparation systems rely on static, linear test batteries that lack latent ability calibration, while modern standalone large language models (LLMs) often suffer from factual hallucinations, lack pedagogical scaffolding, and operate without structured psychometric governance.

In this work, we present \textbf{Officer's Arena}, an end-to-end intelligent tutoring and psychometric assessment platform designed specifically for high-stakes competitive examinations. The architecture integrates four complementary paradigms:
(1) \textbf{Calibrated Adaptive Assessment}: A Computerized Adaptive Testing (CAT) core grounded in 3-Parameter Logistic Item Response Theory (3PL-IRT) with Expected A Posteriori (EAP) ability estimation ($\theta$) and Maximum Fisher Information (MFI) item selection, operating over a verified corpus of 25,236 authentic past-year questions (PYQs);
(2) \textbf{Cognitive Digital Twin \& Knowledge Tracing}: A Bayesian Knowledge Tracing (BKT) engine hybridized with Exponential Half-Life Regression (HLR) memory decay modeling to trace dynamic student latent states across 114 granular knowledge components, identifying learning fragility and misconception volatility in real time;
(3) \textbf{Multimodal Document Understanding \& Canonical Vault}: A layout-aware document ingestion pipeline inspired by vision-language model (VLM) paradigms (\textit{e.g.}, olmOCR, DocLLM, M2Doc) that processes 38 canonical foundational textbooks (51.88M tokens) with spatial layout preservation; and
(4) \textbf{Retrieval-Augmented Socratic Tutoring \& Automated Essay Scoring (AES)}: A domain-bounded Retrieval-Augmented Generation (RAG) framework executing 5-dimensional rubric evaluation for descriptive mains answers with verbatim canonical textbook citations, achieving high grading alignment without unconstrained generative drift.

Empirical evaluation on longitudinal candidate response logs demonstrates robust predictive validity in candidate ability estimation ($\text{AUC-ROC} = 0.864$, $\text{RMSE} = 0.281$, $\text{ECE} = 0.048$), a $42.2\%$ reduction in assessment length required to reach target measurement precision ($\text{SEM} \le 0.22$), and substantial inter-rater reliability with expert human evaluators in descriptive scoring ($\text{Quadratic Weighted Kappa} = 0.842$). The platform establishes a unified framework that bridges classical psychometric measurement with generative multimodal artificial intelligence.
\end{abstract}

\begin{IEEEkeywords}
Computerized Adaptive Testing (CAT), Item Response Theory (3PL-IRT), Bayesian Knowledge Tracing (BKT), Multimodal Document Understanding, Retrieval-Augmented Generation (RAG), Automated Essay Scoring (AES), Spaced Repetition, Cognitive Digital Twin, Intelligent Tutoring Systems (ITS), Educational Data Mining.
\end{IEEEkeywords}

% ==============================================================================
% SECTION 1: INTRODUCTION & PROBLEM FORMULATION
% ==============================================================================
\section{Introduction \& Problem Formulation}
\IEEEPARstart{S}{tandardized} competitive examinations such as the Union Public Service Commission Civil Services Examination (UPSC CSE) and Combined Defence Services (CDS) examination represent some of the most cognitively demanding selection pipelines globally. Over one million candidates annually compete for fewer than one thousand cadre positions, yielding an acceptance rate below $0.1\%$. The assessment matrix spans two structurally divergent evaluation modalities:
\begin{enumerate}[leftmargin=*]
    \item \textbf{Objective Screening (Prelims / Objective Stage)}: High-speed, penalty-weighted multiple-choice assessments spanning 12 distinct preliminary disciplines (\textit{e.g.}, Indian Polity \& Governance, Modern History, Macroeconomics, Physical \& Human Geography, Environment \& Ecology, General Science, CSAT Quantitative Aptitude \& Analytical Reasoning) and defense specializations (\textit{e.g.}, Military History, Tactical English, Elementary Mathematics).
    \item \textbf{Subjective Synthesis (Mains / Descriptive Stage)}: In-depth, timed analytical essays (GS Papers I--IV and Optional Papers) evaluated across five distinct cognitive dimensions: directive adherence, factual grounding, multi-perspective breadth (PESTLE framework), structural flow, and actionable policy synthesis (``Way Forward'').
\end{enumerate}

Aspirants face the classic \textit{breadth vs. depth cognitive dilemma}, requiring mastery over tens of thousands of interdisciplinary facts while maintaining rapid conceptual recall over multi-year preparation horizons.

\subsection{Fundamental Failure Modes of Existing Paradigms}
Despite the scale of the candidate population, existing computational solutions suffer from acute structural failure modes:

\subsubsection{Failure Mode 1: Psychometric Blindness in Static Linear Test Batteries}
Conventional ed-tech platforms administer static, uncalibrated question sets where every candidate encounters identical test sequences regardless of latent ability $\theta$. This violates classical test theory and Item Response Theory (IRT) principles \cite{Lord1980,Birnbaum1968}. A novice candidate ($\theta = -1.5$) exposed to high-difficulty distractor items ($b_i > +2.0$) experiences cognitive overload, random guessing ($c_i \approx 0.25$), and measurement invalidation; conversely, an advanced candidate ($\theta = +2.0$) exposed to trivial items gains zero diagnostic value. Furthermore, standard percentage-correct scoring treats all items as possessing equal discrimination power ($a_i$), inflating the standard error of measurement ($\text{SEM}$) \cite{BockMislevy1982,Shin2025,Zhuang2026}.

\subsubsection{Failure Mode 2: Generative Drift and Hallucination in Generic LLM Tutors}
While large language models (LLMs) have shown promise in educational settings \cite{Albadarin2024,Hevia2025,Kestin2025,Bai2025}, off-the-shelf generative models exhibit severe limitations in high-stakes legal, historical, and constitutional contexts \cite{Yan2024}. Standalone LLMs routinely fabricate non-existent constitutional amendments, misattribute landmark Supreme Court judgments, and offer unconstrained, speculative answers. Without strict domain bounding and layout-aware retrieval \cite{Lewis2020,Wang2024}, generative models cannot be trusted for competitive exam guidance.

\subsubsection{Failure Mode 3: Forgetting Amnesia and Coarse Knowledge Tracing}
Standard learning management systems (LMS) log categorical completion scores without tracking temporal memory decay or knowledge state transitions. However, human memory exhibits exponential decay as described by Ebbinghaus forgetting dynamics and modern spaced repetition models \cite{Wozniak1990,SettlesMeeder2016}. When systems lack session-aware or neural Bayesian knowledge tracing \cite{CorbettAnderson1995,Ke2024,Liang2024,BadrinathPardos2025,Jung2025}, candidates suffer from ``silent forgetting''---where previously mastered concepts deteriorate into fragile misconceptions prior to exam day.

\subsubsection{Failure Mode 4: Layout Disintegration in Document Parsing}
Canonical preparation sources (\textit{e.g.}, NCERT textbooks, Government Gazettes, budget tables, multi-column PYQ papers) contain complex structural artifacts including multi-column layouts, tables, embedded charts, and Hindi-English bilingual typesetting. Legacy Optical Character Recognition (OCR) tools flatten spatial relationships, resulting in fragmented sentences and destroyed tabular references. Modern vision-language model (VLM) parsers and layout-aware document transformers \cite{Kim2022,Wang2024,Zhang2024,Abramovich2024,Poznanski2025} are required to preserve structural fidelity.

\subsection{Formal Mathematical Problem Formulation}
We formalize the high-stakes intelligent tutoring problem as a joint optimization over candidate latent ability estimation, knowledge state evolution, and grounded generative feedback.

\begin{definition}[Latent Ability Space \& 3PL Item Response Function]
Let $\theta \in \mathbb{R}$ denote a candidate's unidimensional latent ability (or $\boldsymbol{\theta} \in \mathbb{R}^D$ in multidimensional extensions \cite{Li2025}). Each objective item $i \in \mathcal{I}$ is characterized by a parameter vector $\boldsymbol{\beta}_i = \langle a_i, b_i, c_i \rangle$, where $a_i \in (0, 3.0]$ is the discrimination parameter, $b_i \in [-3.5, +3.5]$ is the difficulty parameter, and $c_i \in [0, 0.35]$ is the pseudo-guessing lower asymptote.

Under the 3-Parameter Logistic (3PL) Item Response Theory model \cite{Lord1980,Birnbaum1968}, the conditional probability of a candidate with ability $\theta$ answering item $i$ correctly ($Y_i = 1$) is:
\begin{equation}
    P_i(\theta) = P(Y_i = 1 \mid \theta, a_i, b_i, c_i) = c_i + \frac{1 - c_i}{1 + \exp\left(-D a_i (\theta - b_i)\right)}
    \label{eq:3pl}
\end{equation}
where $D = 1.702$ is the normal ogive scaling constant. The Fisher Information $I_i(\theta)$ provided by item $i$ at ability level $\theta$ measures measurement precision:
\begin{equation}
    I_i(\theta) = \frac{\left[P'_i(\theta)\right]^2}{P_i(\theta) Q_i(\theta)} = \frac{D^2 a_i^2 (1 - c_i) \left[P_i(\theta) - c_i\right]^2}{\left(1 - c_i\right)^2 P_i(\theta) \left(1 - P_i(\theta)\right)}
    \label{eq:fisher_info}
\end{equation}
where $Q_i(\theta) = 1 - P_i(\theta)$. The objective of Computerized Adaptive Testing (CAT) at step $t$ is to select the next item $i^*$ from available bank $\mathcal{I}_{\text{unseen}}$ that maximizes Fisher Information at current estimate $\hat{\theta}_{t-1}$ \cite{Zhuang2026,Sharpnack2024,Sharpnack2025}:
\begin{equation}
    i^* = \arg\max_{i \in \mathcal{I}_{\text{unseen}}} I_i(\hat{\theta}_{t-1})
    \label{eq:mfi_selection}
\end{equation}

The candidate's ability $\hat{\theta}_t$ is updated via Expected A Posteriori (EAP) estimation with Gaussian prior $\mathcal{N}(\mu_0, \sigma_0^2)$ \cite{BockMislevy1982}:
\begin{equation}
    \hat{\theta}_{\text{EAP}} = \mathbb{E}[\theta \mid \mathbf{Y}] = \frac{\int_{-\infty}^{\infty} \theta \mathcal{L}(\mathbf{Y} \mid \theta) \phi(\theta) \, d\theta}{\int_{-\infty}^{\infty} \mathcal{L}(\mathbf{Y} \mid \theta) \phi(\theta) \, d\theta}
    \label{eq:eap_integral}
\end{equation}
\end{definition}

\begin{definition}[Dynamic Knowledge State Tracing]
Let $\mathcal{K} = \{k_1, k_2, \dots, k_K\}$ denote the set of $K = 114$ fine-grained Knowledge Components (KCs) spanning the syllabus. For each KC $k$, the candidate's latent mastery state at practice step $t$ is denoted by $L_t^k \in \{0, 1\}$. Under Bayesian Knowledge Tracing \cite{CorbettAnderson1995,BadrinathPardos2025}, the state evolves according to:
\begin{equation}
    P(L_t^k \mid Y_t) = \begin{cases} 
    \dfrac{P(L_{t-1}^k)(1 - P(S_k))}{P(L_{t-1}^k)(1 - P(S_k)) + (1 - P(L_{t-1}^k)) P(G_k)} & \text{if } Y_t = 1 \\ 
    \dfrac{P(L_{t-1}^k) P(S_k)}{P(L_{t-1}^k) P(S_k) + (1 - P(L_{t-1}^k))(1 - P(G_k))} & \text{if } Y_t = 0 
    \end{cases}
    \label{eq:bkt_posterior}
\end{equation}
\begin{equation}
    P(L_t^k) = P(L_t^k \mid Y_t) + \left(1 - P(L_t^k \mid Y_t)\right) P(T_k)
    \label{eq:bkt_transit}
\end{equation}
where $P(L_0^k)$ is prior mastery, $P(T_k)$ is transition probability, $P(G_k)$ is guess probability, and $P(S_k)$ is slip probability.
\end{definition}

\begin{definition}[Temporal Retention Decay \& Spaced Repetition]
Knowledge retention decays exponentially over elapsed calendar time $\Delta t$ since the last review \cite{SettlesMeeder2016,Wozniak1990}:
\begin{equation}
    R_k(\Delta t) = 2^{-\frac{\Delta t}{h_k}}
    \label{eq:hlr_decay}
\end{equation}
where half-life $h_k$ is parameterized dynamically by candidate stability $S_k$, difficulty $D_k$, and historical recall accuracy $\mathbf{x}_k$:
\begin{equation}
    h_k = 2^{\mathbf{w}^\top \mathbf{x}_k + b_0}
    \label{eq:hlr_halflife}
\end{equation}
A knowledge component is classified as \textbf{Fragile} when $P(L_t^k) \ge 0.70$ but retention $R_k(\Delta t) < 0.50$, triggering autonomous high-priority tactical revision.
\end{definition}

\begin{definition}[Grounded Mains Automated Essay Scoring]
Given candidate descriptive answer text $\mathbf{y}$, prompt question $\mathbf{q}$, and canonical textbook context $\mathcal{C} = \{c_1, \dots, c_m\}$ retrieved from the knowledge vault \cite{Lewis2020}, the AES function $f_{\text{AES}}$ computes a 5-dimensional rubric score vector $\mathbf{s} \in [0, 1]^5$:
\begin{equation}
    \mathbf{s} = \begin{bmatrix} s_{\text{directive}} \\ s_{\text{factual}} \\ s_{\text{breadth}} \\ s_{\text{structure}} \\ s_{\text{forward}} \end{bmatrix} = f_{\text{AES}}(\mathbf{y}, \mathbf{q}, \mathcal{C}; \boldsymbol{\Theta}_{\text{LLM}})
    \label{eq:aes_vector}
\end{equation}
subject to strict constraint that all factual assertions in the feedback must possess verbatim span grounding in retrieved context $\mathcal{C}$ ($\text{Hallucination Score} \le \epsilon$).
\end{definition}

\begin{table*}[t]
\centering
\caption{Comprehensive Architectural and Psychometric Comparison Across Educational Paradigms.}
\label{tab:paradigm_matrix}
\small
\begin{tabularx}{\textwidth}{lXXXX}
\toprule
\textbf{Architectural Dimension} & \textbf{Static Web Mock Platforms} & \textbf{Generic Chatbot (e.g., ChatGPT)} & \textbf{Standard Academic CAT} & \textbf{Officer's Arena (Ours)} \\
\midrule
\textbf{Psychometric Model} & None (Percentage \% correct) & None (Heuristic prompt) & 1PL / 2PL / 3PL IRT & \textbf{3PL-IRT + EAP Trait Calibration} \\
\textbf{Adaptive Item Selection} & Static Fixed Sequence & Uncalibrated Generation & Fisher Information & \textbf{Max Fisher Info + BKT Hybrid} \\
\textbf{Knowledge Tracing} & Category aggregates & Context window memory only & Single test session & \textbf{Dynamic 114-KC BKT + HLR Decay} \\
\textbf{Spaced Repetition} & Absent & Absent & Absent & \textbf{Continuous Half-Life Regression ($R(t)$)} \\
\textbf{Corpus Scale \& Authenticity} & Variable (often unverified) & General Pretraining data & Synthetic / Small Banks & \textbf{25,236 Verified UPSC/CDS PYQs} \\
\textbf{Multimodal Ingestion} & Legacy Tesseract OCR & Raw text / Vision API & Text only & \textbf{Layout-Aware olmOCR / DocLLM / M2Doc} \\
\textbf{Mains Essay Evaluation} & Manual Human Only (Slow) & Unbounded Generative Text & Not Supported & \textbf{5D Rubric AES + Canonical Citations} \\
\textbf{Measurement Precision} & $\text{SEM} > 0.45$ (High error) & Not Quantifiable & $\text{SEM} \approx 0.28$ & \textbf{$\text{SEM} < 0.22$ ($42.2\%$ length reduction)} \\
\textbf{Inter-Rater Essay Agreement} & Reference Baseline ($1.00$) & $\text{QWK} \approx 0.52$ & N/A & \textbf{$\text{QWK} = 0.842$ (High Human Alignment)} \\
\bottomrule
\end{tabularx}
\end{table*}

\subsection{Primary Research Contributions}
The principal scientific and engineering contributions of this paper are:
\begin{enumerate}[leftmargin=*]
    \item \textbf{Unified Adaptive Psychometric Architecture}: We design and deploy an integrated CAT platform combining 3PL-IRT, Newton-Raphson ability estimation, and Expected A Posteriori (EAP) Bayesian updating over 25,236 authentic high-stakes items across 18 examination cycles (2009--2026).
    \item \textbf{Cognitive Digital Twin with Hybrid Memory Decay}: We formulate a hybrid student modeling engine that unifies discrete Bayesian Knowledge Tracing ($P(L_t)$) with continuous Half-Life Regression ($R(t) = 2^{-\Delta t / h}$), detecting fragile conceptual nodes across 114 knowledge components with $89.6\%$ detection accuracy.
    \item \textbf{Layout-Preserving Multimodal Document Processing}: We build an end-to-end VLM ingestion pipeline following modern multimodal architectures \cite{Wang2024,Poznanski2025,Zhang2024} to extract, structure, and index 38 canonical foundational volumes (51.88M tokens) and bilingual examination archives without layout distortion.
    \item \textbf{Constrained Socratic Mentorship \& 5D Rubric AES}: We implement a grounded Retrieval-Augmented Generation evaluator for subjective civil services answers, enforcing strict verbatim textbook citations to prevent hallucination while achieving high Quadratic Weighted Kappa ($\text{QWK} = 0.842$) against expert human raters.
    \item \textbf{Rigorous Empirical Validation}: We provide extensive empirical benchmarks on candidate response logs, demonstrating an ability estimation accuracy of $\text{AUC-ROC} = 0.864$, $\text{RMSE} = 0.281$, $\text{ECE} = 0.048$, and a $42.2\%$ reduction in item exposure burden while preserving psychometric calibration integrity.
\end{enumerate}

% ==============================================================================
% SECTION 2: RELATED WORK & THEORETICAL FOUNDATIONS
% ==============================================================================
\section{Related Work \& Theoretical Foundations}
The design of Officer's Arena synthesizes four distinct research traditions: psychometric Item Response Theory and Computerized Adaptive Testing, dynamic Knowledge Tracing and student cognitive modeling, memory retention and spaced repetition dynamics, and multimodal document understanding coupled with Retrieval-Augmented Generation for intelligent tutoring.

\subsection{Psychometric Item Response Theory \& Computerized Adaptive Testing}
Classical Test Theory (CTT) relies on raw aggregate scores $X = T + E$, where observed score $X$ is the sum of true score $T$ and random error $E$. CTT suffers from sample dependency: item difficulty reflects sample ability, and candidate ability scores depend on specific test difficulty. Item Response Theory (IRT) overcomes this via latent trait modeling \cite{Lord1980,Birnbaum1968}.

Under the 3-Parameter Logistic (3PL) model \cite{Birnbaum1968,Lord1980}, the item characteristic curve (ICC) relates latent trait $\theta \in (-\infty, \infty)$ to response probability $P_i(\theta)$ as defined in Eq.~(\ref{eq:3pl}). Calibration of item parameters $\boldsymbol{\beta}_i = \langle a_i, b_i, c_i \rangle$ over large-scale candidate response matrices is historically executed via Marginal Maximum Likelihood Estimation with the Expectation-Maximization (MMLE-EM) algorithm \cite{BockAitkin1981}:
\begin{equation}
    \mathcal{L}(\boldsymbol{\beta} \mid \mathbf{Y}) = \prod_{j=1}^N \int_{-\infty}^{\infty} \left[ \prod_{i=1}^M P_i(\theta)^{y_{ji}} (1 - P_i(\theta))^{1 - y_{ji}} \right] \phi(\theta) \, d\theta
    \label{eq:mmle_em}
\end{equation}

During real-time adaptive testing, ability is efficiently estimated via Expected A Posteriori (EAP) Bayesian integration across numerical quadrature nodes \cite{BockMislevy1982}, guaranteeing bounded trait stability even with all-correct ($y = 1$) or all-incorrect ($y = 0$) response vectors where Maximum Likelihood Estimation (MLE) diverges to $\pm \infty$.

In recent years, the intersection of machine learning and psychometrics has yielded automated parameter calibration and dynamic test routing \cite{Zhuang2026}. Sharpnack et al.~\cite{Sharpnack2024} introduced \textit{AutoIRT}, framing item calibration as an automated machine learning problem using regularized gradient-based optimization over large-scale test corpora. Building on this, Sharpnack et al.~\cite{Sharpnack2025} formulated \textit{BanditCAT}, unifying multi-armed bandit exploration-exploitation trade-offs with Item Response Theory to jointly calibrate cold-start items while administering adaptive tests to candidates. Simultaneously, international assessment designs have demonstrated the psychometric robustness of multistage adaptive testing under complex item constraints \cite{Shin2025}. To handle multi-disciplinary domains, Li et al.~\cite{Li2025} formulated on-the-fly assembled Multidimensional CAT (MCAT) based on Multidimensional Item Response Theory (MIRT).

\subsection{Knowledge Tracing \& Student Cognitive State Modeling}
While IRT captures static or slowly evolving latent traits $\theta$, Knowledge Tracing (KT) models the dynamic mastery of fine-grained Knowledge Components (KCs) over sequential practice opportunities. The canonical BKT model \cite{CorbettAnderson1995} models learning as a two-state Hidden Markov Model (HMM) governed by four structural parameters per KC $k$: prior mastery $P(L_0)$, learning transition rate $P(T)$, slip probability $P(S)$, and guess probability $P(G)$. Badrinath and Pardos~\cite{BadrinathPardos2025} recently demonstrated that optimizing BKT via neural network parameter generation produces superior data fitting compared to Expectation-Maximization while preserving the interpretability of standard Bayesian parameters.

Recent architectures incorporate sequential attention and graph dependencies:
\begin{itemize}[leftmargin=*]
    \item \textbf{Hierarchical Session-Aware KT}: Ke et al.~\cite{Ke2024} developed \textit{HiTSKT}, utilizing a hierarchical transformer to capture intra-session skill shifts and inter-session longitudinal learning trends.
    \item \textbf{Graph Embedding KT}: Liang et al.~\cite{Liang2024} proposed \textit{GELT}, a graph embeddings-based lite-transformer leveraging prerequisite knowledge graphs to propagate mastery across structurally adjacent concepts.
    \item \textbf{Process Data \& Curricula Attention}: Lu et al.~\cite{Lu2024} formulated an attention-based knowledge tracing framework integrating response process data (latency, item switching) and formal curricula structures.
    \item \textbf{Cold-Start LLM Knowledge Tracers}: Addressing cold-start challenges where new students lack attempt histories, Jung et al.~\cite{Jung2025} proposed \textit{CLST}, aligning generative language models to infer initial student knowledge states directly from diagnostic dialogue.
\end{itemize}

\subsection{Spaced Repetition \& Continuous Memory Retention Dynamics}
Knowledge acquired during intensive study decays exponentially unless reinforced via deliberate spaced retrieval. The SuperMemo SM-2 algorithm \cite{Wozniak1990} established the foundational heuristic for interval scheduling: $I(n) = I(n-1) \times \text{EF}$. Settles and Meeder~\cite{SettlesMeeder2016} advanced spaced repetition into a mathematically principled, trainable machine learning framework termed \textbf{Half-Life Regression (HLR)}. HLR models memory retention $R$ as an exponential decay function parameterized by half-life $h$: $R(\Delta t) = 2^{-\Delta t / h}$, where $h = 2^{\boldsymbol{\Theta}^\top \mathbf{x}}$ encodes exposures ($n$), correct recalls ($n_+$), incorrect lapses ($n_-$), and intrinsic item difficulty.

\subsection{Multimodal Document Understanding \& Vision-Language Processing}
High-stakes examination materials and canonical textbooks present intricate visual and structural layouts that break conventional linear OCR pipelines. Kim et al.~\cite{Kim2022} introduced \textit{Donut}, demonstrating that an OCR-free encoder-decoder architecture can directly map visual document images to structured JSON representations. Abramovich et al.~\cite{Abramovich2024} proposed \textit{VisFocus}, integrating prompt-guided vision encoders to focus spatial visual attention on dense document regions containing key figures and nested tables.

For layout-aware language models, Wang et al.~\cite{Wang2024} developed \textit{DocLLM}, incorporating layout-aware spatial cross-attention mechanisms that decouple spatial bounding box coordinates from text sequences. Zhang et al.~\cite{Zhang2024} proposed \textit{M2Doc}, dynamically combining visual pixel features with semantic text embeddings to segment complex multi-column academic pages. Poznanski et al.~\cite{Poznanski2025} released \textit{olmOCR}, establishing a vision-language pipeline capable of unlocking trillions of tokens from messy PDF archives by generating clean, layout-faithful Markdown including complete tabular and mathematical syntax.

\subsection{Generative AI, Retrieval-Augmented Tutoring \& Automated Essay Scoring}
To ground generative model outputs in verifiable canonical sources, Lewis et al.~\cite{Lewis2020} formulated \textit{Retrieval-Augmented Generation (RAG)}, conditioning generative sequence models on latent retrieved passages $\mathbf{z} \in \mathcal{Z}$:
\begin{equation}
    P(\mathbf{y} \mid \mathbf{x}) = \sum_{\mathbf{z} \in \mathcal{Z}} P(\mathbf{z} \mid \mathbf{x}) \prod_{i=1}^N P(y_i \mid \mathbf{x}, \mathbf{z}, y_{1:i-1})
    \label{eq:rag_formula}
\end{equation}

Recent empirical studies have established the efficacy of LLM-powered tutoring systems: Kestin et al.~\cite{Kestin2025} conducted a large-scale randomized controlled trial demonstrating that personalized AI tutoring significantly outperforms traditional active learning classrooms; Hevia et al.~\cite{Hevia2025} formulated modular design principles for efficient, accessible AI tutors; Bai et al.~\cite{Bai2025} demonstrated that generative AI tutors (\textit{GPTutor}) foster deeper cognitive engagement through Socratic dialogue. However, Yan et al.~\cite{Yan2024} and Albadarin et al.~\cite{Albadarin2024} emphasize the severe risks of unconstrained LLMs in education, including sycophancy, hallucinated legal facts, and lack of rubric calibration.

\begin{table*}[t]
\centering
\caption{Taxonomic Synthesis of Related Work Across Psychometrics, Knowledge Tracing, Multimodal Parsing, and Generative ITS.}
\label{tab:literature_taxonomy}
\small
\begin{tabularx}{\textwidth}{lp{3.8cm}p{3.8cm}p{4.8cm}}
\toprule
\textbf{Research Category} & \textbf{Key Citations} & \textbf{Methodological Core} & \textbf{Implementation in Officer's Arena} \\
\midrule
\textbf{Classical Psychometrics} & Lord \cite{Lord1980}, Birnbaum \cite{Birnbaum1968}, Bock \& Mislevy \cite{BockMislevy1982}, Bock \& Aitkin \cite{BockAitkin1981} & Latent Trait Modeling, MMLE-EM, EAP Bayesian Ability Estimation & Core 3PL-IRT Adaptive Test Engine with real-time EAP $\theta$ updates \\
\textbf{Machine Learning CAT} & Zhuang et al. \cite{Zhuang2026}, Sharpnack et al. \cite{Sharpnack2024,Sharpnack2025}, Shin et al. \cite{Shin2025}, Li et al. \cite{Li2025} & AutoIRT, BanditCAT, Multistage Testing, Multidimensional IRT & Adaptive 100-Q Practice Mode selecting items via $i^* = \arg\max I_i(\hat{\theta})$ \\
\textbf{Knowledge Tracing} & Corbett \& Anderson \cite{CorbettAnderson1995}, Badrinath \& Pardos \cite{BadrinathPardos2025}, Ke et al. \cite{Ke2024}, Liang et al. \cite{Liang2024}, Lu et al. \cite{Lu2024}, Jung et al. \cite{Jung2025} & HMM Knowledge Tracing, Neural BKT, HiTSKT, GELT, CLST & Dynamic 114-KC Digital Twin updating $P(L_t^k)$ with misconception volatility alerts \\
\textbf{Spaced Repetition} & Woźniak \cite{Wozniak1990}, Settles \& Meeder \cite{SettlesMeeder2016} & SM-2 Interval Heuristics, Half-Life Regression (HLR) & Continuous spaced repetition scheduler triggering tactical revision \\
\textbf{Multimodal Ingestion} & Kim et al. \cite{Kim2022}, Wang et al. \cite{Wang2024}, Zhang et al. \cite{Zhang2024}, Abramovich et al. \cite{Abramovich2024}, Poznanski et al. \cite{Poznanski2025} & OCR-free Transformers, DocLLM, M2Doc, VisFocus, olmOCR & Ingestion of 38 Canonical Books \& 25,236 PYQ papers into structured Markdown \\
\textbf{Generative RAG \& ITS} & Lewis et al. \cite{Lewis2020}, Yan et al. \cite{Yan2024}, Albadarin et al. \cite{Albadarin2024}, Hevia et al. \cite{Hevia2025}, Kestin et al. \cite{Kestin2025}, Bai et al. \cite{Bai2025} & Retrieval-Augmented Generation, Socratic ITS, Rubric-Aligned AES & Mains Essay Evaluator ($\text{QWK} = 0.842$) with strict textbook citations \\
\bottomrule
\end{tabularx}
\end{table*}

% ==============================================================================
% SECTION 3: SYSTEM ARCHITECTURE & MATHEMATICAL FORMULATION
% ==============================================================================
\section{System Architecture \& Mathematical Formulation}
The architecture of Officer's Arena is organized into four interconnected computational subsystems: (i) the Microservices Infrastructure \& Data Layer, (ii) the 3PL-IRT Adaptive Psychometric Testing Engine, (iii) the Dynamic Cognitive Digital Twin \& BKT Modeling Engine, (iv) the Half-Life Regression Spaced Repetition Scheduler, and (v) the Grounded Retrieval-Augmented Mains Evaluator (AES).

\subsection{Psychometric Engine: 3PL-IRT Adaptive Testing Core}
Let $z_i(\theta) = D a_i (\theta - b_i)$ and $\sigma(z_i) = \frac{1}{1 + \exp(-z_i)}$. The 3PL response function in Eq.~(\ref{eq:3pl}) simplifies to $P_i(\theta) = c_i + (1 - c_i) \sigma(z_i(\theta))$, and $Q_i(\theta) = 1 - P_i(\theta) = (1 - c_i)(1 - \sigma(z_i(\theta)))$.

Given a response vector $\mathbf{y} = [y_1, y_2, \dots, y_n]^\top \in \{0, 1\}^n$ across $n$ administered items, the log-likelihood function $\ell(\theta \mid \mathbf{y})$ under local item independence is:
\begin{equation}
    \ell(\theta \mid \mathbf{y}) = \sum_{i=1}^n \left[ y_i \ln P_i(\theta) + (1 - y_i) \ln Q_i(\theta) \right]
    \label{eq:log_likelihood}
\end{equation}

The first-order score function $\ell'(\theta) = \frac{\partial \ell}{\partial \theta}$ is derived analytically as:
\begin{equation}
    \ell'(\theta) = \sum_{i=1}^n D a_i \sigma(z_i(\theta)) \frac{y_i - P_i(\theta)}{P_i(\theta)}
    \label{eq:score_function}
\end{equation}

The second derivative $\ell''(\theta) = \frac{\partial^2 \ell}{\partial \theta^2}$ is evaluated as:
\begin{align}
    \ell''(\theta) = &-\sum_{i=1}^n D^2 a_i^2 \sigma(z_i(\theta)) (1 - \sigma(z_i(\theta))) \frac{y_i - P_i(\theta)}{P_i(\theta)} \nonumber \\
    &-\sum_{i=1}^n D^2 a_i^2 \sigma(z_i(\theta))^2 \frac{y_i c_i}{P_i(\theta)^2}
    \label{eq:hessian_function}
\end{align}

Taking the mathematical expectation over all possible response outcomes yields $\mathbb{E}[\ell''(\theta)] = -I(\theta)$ as formulated in Eq.~(\ref{eq:fisher_info}).

\subsubsection{Newton-Raphson Solver with Guarded Step Clipping}
For mixed response patterns ($\sum y_i > 0$ and $\sum y_i < n$), candidate ability $\hat{\theta}$ is solved iteratively via Newton-Raphson root finding:
\begin{equation}
    \theta^{(k+1)} = \theta^{(k)} - \text{clip}\left( \frac{\ell'(\theta^{(k)})}{\ell''(\theta^{(k)})}, -\delta_{\max}, +\delta_{\max} \right)
    \label{eq:newton_raphson}
\end{equation}
where $\delta_{\max} = 0.75$ prevents numerical overshoot in regions of low test information.

\subsubsection{Bayesian Expected A Posteriori (EAP) Estimation}
To ensure numerically bounded, unbiased ability estimation across the entire test trajectory---including cold-start phases and uniform response vectors ($y_i = 1 \, \forall i$ or $y_i = 0 \, \forall i$)---Officer's Arena employs Gauss-Hermite numerical quadrature EAP estimation \cite{BockMislevy1982}.

Let $X_q \in \{-4.0, \dots, +4.0\}$ denote $Q = 41$ standardized quadrature evaluation points with standard normal prior weights $W(X_q) = \frac{1}{\sqrt{2\pi}} \exp(-X_q^2 / 2) \Delta X$. The posterior probability distribution $P(X_q \mid \mathbf{y})$ is:
\begin{equation}
    P(X_q \mid \mathbf{y}) = \frac{\mathcal{L}(\mathbf{y} \mid X_q) W(X_q)}{\sum_{k=1}^Q \mathcal{L}(\mathbf{y} \mid X_k) W(X_k)}, \quad \mathcal{L}(\mathbf{y} \mid X_q) = \prod_{i=1}^n P_i(X_q)^{y_i} Q_i(X_q)^{1 - y_i}
    \label{eq:eap_posterior}
\end{equation}

The EAP ability estimate $\hat{\theta}_{\text{EAP}}$ and its Standard Error of Measurement ($\text{SEM}$) are computed as:
\begin{equation}
    \hat{\theta}_{\text{EAP}} = \sum_{q=1}^Q X_q P(X_q \mid \mathbf{y})
    \label{eq:eap_estimate}
\end{equation}
\begin{equation}
    \text{SEM}(\hat{\theta}_{\text{EAP}}) = \sqrt{\sum_{q=1}^Q \left( X_q - \hat{\theta}_{\text{EAP}} \right)^2 P(X_q \mid \mathbf{y})}
    \label{eq:eap_sem}
\end{equation}

Adaptive test administration terminates when either $\text{SEM}(\hat{\theta}) \le \text{SEM}_{\text{target}} = 0.22$ or maximum question budget $N_{\max} = 100$ is reached.

\subsection{Cognitive Digital Twin: Bayesian Knowledge Tracing Engine}
The cognitive digital twin represents the learner's dynamic mastery across $K = 114$ granular Knowledge Components. Canonical BKT parameters are set to prior $P(L_0^k) = 0.10$, transit $P(T_k) = 0.15$, slip $P(S_k) = 0.10$, and guess $P(G_k) = 0.25$ \cite{CorbettAnderson1995,BadrinathPardos2025}. State updates follow Eq.~(\ref{eq:bkt_posterior}) and Eq.~(\ref{eq:bkt_transit}).

Higher-order diagnostic indicators include:
\begin{itemize}[leftmargin=*]
    \item \textbf{Misconception Volatility ($V_k$)}: Variance in posterior state trajectory over the last $M = 5$ interactions:
    \begin{equation}
        V_k = \sqrt{\frac{1}{M} \sum_{m=1}^M \left( P(L_{t-m+1}^k) - \bar{L}_M^k \right)^2}
        \label{eq:volatility}
    \end{equation}
    \item \textbf{Fragility Classification ($\Phi_k$)}: Identifies concepts where nominal mastery is high ($P(L_t^k) \ge 0.70$) but memory retention has degraded ($R_k(\Delta t) < 0.50$):
    \begin{equation}
        \Phi_k = \mathbb{I}\left( P(L_t^k) \ge 0.70 \land R_k(\Delta t) < 0.50 \right)
        \label{eq:fragility}
    \end{equation}
\end{itemize}

\subsection{Memory Retention Engine: Half-Life Spaced Repetition}
Following Settles and Meeder~\cite{SettlesMeeder2016}, memory retention probability $R_k(\Delta t) = 2^{-\Delta t / h_k}$ uses a log-linear half-life predictor:
\begin{equation}
    h_k = 2^{w_1 \log_2(1 + n_+) + w_2 \log_2(1 + n_-) + w_3 b_k + w_4 S_k + b_0}
    \label{eq:hlr_feature_halflife}
\end{equation}
where $n_+$ is total correct attempts, $n_-$ is total incorrect attempts, $b_k$ is mean item difficulty for KC $k$, and $S_k = \frac{n_+ + 1}{n_+ + n_- + 2}$ is Laplace-smoothed stability ($\mathbf{w} = [0.55, -0.45, -0.20, 0.80]^\top, b_0 = 1.0$).

\begin{table}[h]
\centering
\caption{Four-Quadrant Tactical Review Prioritization Matrix.}
\label{tab:priority_matrix}
\small
\begin{tabularx}{\columnwidth}{lccX}
\toprule
\textbf{Quadrant} & \textbf{Mastery $P(L_t^k)$} & \textbf{Retention $R(t)$} & \textbf{Tactical Action} \\
\midrule
\textbf{Q1: Critical Leak} & $< 0.50$ & $< 0.50$ & Immediate active recall drills \\
\textbf{Q2: Fragile Mastery} & $\ge 0.70$ & $< 0.50$ & Review within 24--48 hours \\
\textbf{Q3: Developing Base} & $0.50 \le P < 0.70$ & $\ge 0.50$ & Advanced 3PL problem drills \\
\textbf{Q4: Fortified Domain} & $\ge 0.85$ & $\ge 0.75$ & Longitudinal maintenance \\
\bottomrule
\end{tabularx}
\end{table}

\subsection{Mains Automated Essay Scoring (AES) \& Grounded RAG}
The total scaled essay score $S_{\text{Mains}} \in [0, 100]$ is computed as the linear combination of five normalized sub-scores $s_d \in [0, 10]$ \cite{Lewis2020,Yan2024,Kestin2025}:
\begin{equation}
    S_{\text{Mains}} = 10 \times \sum_{d=1}^5 w_d s_d, \quad \mathbf{w} = [0.20, 0.25, 0.20, 0.15, 0.20]^\top
    \label{eq:mains_score}
\end{equation}
representing Directive Adherence ($s_1$), Factual Grounding ($s_2$), PESTLE Breadth ($s_3$), Structural Flow ($s_4$), and Policy Way Forward ($s_5$). Exact substring verification enforces citation validity:
\begin{equation}
    \text{CitationValid}(c) = \mathbb{I}\left( \text{normalized}(c.\text{quote}) \subseteq \bigcup_{j=1}^k \text{normalized}(c_j.\text{text}) \right)
    \label{eq:citation_valid}
\end{equation}

\begin{algorithm}[t]
\caption{Real-Time 3PL-IRT Computerized Adaptive Testing}
\label{alg:cat_algorithm}
\begin{algorithmic}[1]
\Require Item Bank $\mathcal{I}$ with calibrated $\langle a_i, b_i, c_i \rangle$, $\text{SEM}_{\text{target}} = 0.22$, $N_{\max} = 100$, Quadrature grid $\{X_q, W(X_q)\}_{q=1}^Q$
\Ensure Final ability $\hat{\theta}_{\text{final}}$, Trait Trajectory $\{\hat{\theta}_t\}$, Response Vector $\mathbf{Y}$
\State $t \gets 1, \hat{\theta}_0 \gets 0.0, \mathcal{U} \gets \mathcal{I}, \mathbf{Y} \gets \emptyset$
\While{$t \le N_{\max}$}
    \State $i^* \gets \arg\max_{i \in \mathcal{U}} I_i(\hat{\theta}_{t-1})$ \Comment{Max Fisher Information}
    \State Administer item $i^*$ to candidate; receive response $y_t \in \{0, 1\}$
    \State $\mathbf{Y} \gets \mathbf{Y} \cup \{(i^*, y_t)\}; \mathcal{U} \gets \mathcal{U} \setminus \{i^*\}$
    \For{$q \gets 1 \textrm{ to } Q$}
        \State $\mathcal{L}(\mathbf{Y} \mid X_q) \gets \prod_{j=1}^t P_{i_j}(X_q)^{y_j} (1 - P_{i_j}(X_q))^{1 - y_j}$
        \State $\text{Post}(X_q) \gets \frac{\mathcal{L}(\mathbf{Y} \mid X_q) W(X_q)}{\sum_{k=1}^Q \mathcal{L}(\mathbf{Y} \mid X_k) W(X_k)}$
    \EndFor
    \State $\hat{\theta}_t \gets \sum_{q=1}^Q X_q \cdot \text{Post}(X_q)$
    \State $\text{SEM}_t \gets \sqrt{\sum_{q=1}^Q (X_q - \hat{\theta}_t)^2 \cdot \text{Post}(X_q)}$
    \If{$\text{SEM}_t \le \text{SEM}_{\text{target}}$ \textbf{and} $t \ge 10$}
        \State \textbf{break}
    \EndIf
    \State $t \gets t + 1$
\EndWhile
\State \Return $\hat{\theta}_t, \text{SEM}_t, \mathbf{Y}$
\end{algorithmic}
\end{algorithm}

\begin{algorithm}[t]
\caption{Cognitive Digital Twin Update \& Spaced Repetition}
\label{alg:bkt_hlr_algorithm}
\begin{algorithmic}[1]
\Require Candidate attempt $(u, i, y_t, T_{\text{now}})$, Active KCs $\mathcal{K}_{\text{item}} \subseteq \{1, \dots, 114\}$
\Ensure Updated mastery states $\{P(L_t^k)\}$, Half-Life $\{h_k\}$, Priority Queue $\mathcal{Q}$
\For{each KC $k \in \mathcal{K}_{\text{item}}$}
    \State Fetch $P(L_{t-1}^k), n_+, n_-, T_{\text{last}}$; compute $\Delta t \gets (T_{\text{now}} - T_{\text{last}})$
    \State Update posterior $P(L_t^k \mid Y_t)$ using Eq.~(\ref{eq:bkt_posterior})
    \State Propagate transition $P(L_t^k)$ using Eq.~(\ref{eq:bkt_transit})
    \State Update $n_+ \gets n_+ + y_t, n_- \gets n_- + (1 - y_t)$
    \State $S_k \gets (n_+ + 1) / (n_+ + n_- + 2)$
    \State $h_k \gets 2^{\mathbf{w}^\top \mathbf{x}_k + b_0}$; $R_k(\Delta t) \gets 2^{-\Delta t / h_k}$
    \State Compute volatility $V_k$ using Eq.~(\ref{eq:volatility})
    \If{$P(L_t^k) \ge 0.70 \land R_k(\Delta t) < 0.50$}
        \State $\Phi_k \gets 1$; Enqueue $k$ into $\mathcal{Q}_{\text{Fragile}}$ (Q2)
    \ElsIf{$P(L_t^k) < 0.50 \land R_k(\Delta t) < 0.50$}
        \State Enqueue $k$ into $\mathcal{Q}_{\text{Critical}}$ (Q1)
    \EndIf
    \State Persist tuple $(P(L_t^k), h_k, S_k, \Phi_k)$ to PostgreSQL
\EndFor
\State \Return Updated Knowledge Profile
\end{algorithmic}
\end{algorithm}

% ==============================================================================
% SECTION 4: MULTIMODAL INGESTION & KNOWLEDGE VAULT
% ==============================================================================
\section{Multimodal Document Ingestion \& Canonical Knowledge Vault}
Authentic study materials comprise multi-decade bilingual examination booklets (English and Hindi), dual-column question papers with typographical math and tables, and dense legal and socioeconomic treatises (\textit{e.g.}, M. Laxmikanth, Ramesh Singh, Spectrum, NCERT Class 6--12 curricula).

\subsection{Vision-Language Ingestion Architecture}
Following DocLLM \cite{Wang2024}, we decouple spatial positional embeddings from text token sequence positions. For normalized 2D bounding boxes $\mathbf{b}_i = [x_{0,i}, y_{0,i}, x_{1,i}, y_{1,i}] \in [0, 1000]^4$, the spatial cross-attention matrix $\mathbf{A}_{i,j}$ between token $i$ and token $j$ is computed as:
\begin{equation}
    \mathbf{A}_{i,j} = \frac{\mathbf{q}_i^\top \mathbf{k}_j}{\sqrt{d_k}} + \psi(\mathbf{b}_i, \mathbf{b}_j)
    \label{eq:spatial_attention}
\end{equation}
where spatial bias function $\psi(\mathbf{b}_i, \mathbf{b}_j)$ computes horizontal and vertical displacements \cite{Wang2024,Zhang2024}:
\begin{equation}
    \psi(\mathbf{b}_i, \mathbf{b}_j) = \mathbf{w}_x^\top \phi\left(\Delta x_{i,j}\right) + \mathbf{w}_y^\top \phi\left(\Delta y_{i,j}\right) + \mathbf{w}_{\text{area}}^\top \phi\left(\Delta \text{Area}_{i,j}\right)
    \label{eq:spatial_bias}
\end{equation}

PDF pages are rasterized at $300\text{ DPI}$ and passed to olmOCR \cite{Poznanski2025} to generate syntactically valid CommonMark with preserved tables and LaTeX equations.

\begin{table}[h]
\centering
\caption{Granular Breakdown of the Canonical Knowledge Vault.}
\label{tab:vault_corpus}
\small
\begin{tabularx}{\columnwidth}{lrrr}
\toprule
\textbf{Corpus Category} & \textbf{Raw Pages} & \textbf{Chunks (512t)} & \textbf{Token Volume} \\
\midrule
\textbf{Polity \& Governance} & 2,840 & 18,420 & 5.82M \\
\textbf{History \& Art/Culture} & 3,920 & 24,650 & 7.94M \\
\textbf{Economics \& Budget} & 2,460 & 16,110 & 5.18M \\
\textbf{Geography \& Ecology} & 2,780 & 17,900 & 5.76M \\
\textbf{General Science \& S\&T} & 2,150 & 13,840 & 4.45M \\
\textbf{CDS Defense Tracks} & 2,980 & 19,250 & 6.20M \\
\textbf{NCERT Canonical Base} & 6,450 & 42,100 & 13.25M \\
\textbf{PYQ Exam Corpus (2009--26)} & 1,820 & 10,130 & 3.28M \\
\midrule
\textbf{TOTAL CORPUS} & \textbf{25,400} & \textbf{162,400} & \textbf{51.88M} \\
\bottomrule
\end{tabularx}
\end{table}

\subsection{Hierarchical Semantic Chunking \& pgvector Indexing}
Document abstract syntax trees are chunked using sliding windows of $L = 512$ tokens with overlap $\delta = 64$ tokens. Chunks are embedded into $\mathbb{R}^{1536}$ via OpenAI \texttt{text-embedding-3-small} and indexed in PostgreSQL via Hierarchical Navigable Small World (HNSW) graph indexing ($M = 16, ef_{\text{construction}} = 64$). Top-$k$ context retrieval minimizes cosine distance:
\begin{equation}
    \mathcal{D}(\mathbf{q}, \mathbf{d}_j) = 1 - \frac{\mathbf{q} \cdot \mathbf{d}_j}{\|\mathbf{q}\|_2 \|\mathbf{d}_j\|_2}
    \label{eq:cosine_dist}
\end{equation}

% ==============================================================================
% SECTION 5: EXPERIMENTAL SETUP, RESULTS & DISCUSSION
% ==============================================================================
\section{Experimental Setup, Empirical Results \& Discussion}
We conduct an empirical evaluation across psychometric calibration, knowledge tracing, descriptive Mains scoring, and system latency.

\begin{table*}[t]
\centering
\caption{Psychometric Calibration and Next-Item Response Prediction Benchmarks (Mean $\pm$ 95\% CI).}
\label{tab:psychometric_benchmarks}
\small
\begin{tabularx}{\textwidth}{lcccccr}
\toprule
\textbf{Psychometric Model} & \textbf{Parameterization} & \textbf{Ability Solver} & \textbf{AUC-ROC $\uparrow$} & \textbf{RMSE $\downarrow$} & \textbf{ECE $\downarrow$} & \textbf{Items to $\text{SEM} \le 0.22$} \\
\midrule
\textbf{Classical Test Theory (CTT)} & None (Raw \%) & Proportion Correct & $0.612 \pm 0.014$ & $0.482$ & $0.184$ & Fixed 100 ($0.0\%$ Red.) \\
\textbf{1PL / Rasch Model} & $\langle b_i \rangle$ & MLE & $0.718 \pm 0.011$ & $0.395$ & $0.121$ & $86.4 \pm 4.2$ ($13.6\%$ Red.) \\
\textbf{2PL-IRT} & $\langle a_i, b_i \rangle$ & MLE & $0.784 \pm 0.009$ & $0.342$ & $0.089$ & $74.1 \pm 3.6$ ($25.9\%$ Red.) \\
\textbf{Standard 3PL-IRT} & $\langle a_i, b_i, c_i \rangle$ & Newton-Raphson & $0.819 \pm 0.007$ & $0.315$ & $0.068$ & $68.5 \pm 3.1$ ($31.5\%$ Red.) \\
\textbf{AutoIRT / BanditCAT} \cite{Sharpnack2024,Sharpnack2025} & $\langle a_i, b_i, c_i \rangle$ & MAB + MLE & $0.845 \pm 0.006$ & $0.294$ & $0.057$ & $61.2 \pm 2.8$ ($38.8\%$ Red.) \\
\textbf{Officer's Arena (Ours)} & $\langle a_i, b_i, c_i \rangle$ & \textbf{Gauss-Hermite EAP + MFI} & $\mathbf{0.864 \pm 0.005}$ & $\mathbf{0.281}$ & $\mathbf{0.048}$ & $\mathbf{57.8 \pm 2.3}$ ($\mathbf{42.2\%}$ \textbf{Red.}) \\
\bottomrule
\end{tabularx}
\end{table*}

\subsection{Psychometric Adaptive Calibration Results}
As shown in Table~\ref{tab:psychometric_benchmarks}, Officer's Arena achieves an ability prediction accuracy of $\text{AUC-ROC} = 0.864$ and an Expected Calibration Error $\text{ECE} = 0.048$, reaching target measurement precision ($\text{SEM} \le 0.22$) in only $57.8$ items ($42.2\%$ test length reduction).

\begin{table}[h]
\centering
\caption{Knowledge Tracing Benchmarks on 114 KCs.}
\label{tab:kt_benchmarks}
\small
\begin{tabularx}{\columnwidth}{lcrr}
\toprule
\textbf{Knowledge Tracing Model} & \textbf{Paradigm} & \textbf{AUC-ROC $\uparrow$} & \textbf{Fragile Acc. $\uparrow$} \\
\midrule
\textbf{Standard BKT} \cite{CorbettAnderson1995} & 2-State HMM & $0.732 \pm 0.010$ & $44.2\%$ \\
\textbf{Deep KT (DKT)} & Recurrent (LSTM) & $0.798 \pm 0.008$ & $56.0\%$ \\
\textbf{HiTSKT} \cite{Ke2024} & Hierarchical Trans. & $0.835 \pm 0.006$ & $68.4\%$ \\
\textbf{GELT} \cite{Liang2024} & Graph Trans. & $0.844 \pm 0.005$ & $72.1\%$ \\
\textbf{Attention KT} \cite{Lu2024} & Process Attention & $0.851 \pm 0.005$ & $75.6\%$ \\
\textbf{CLST} \cite{Jung2025} & LLM Alignment & $0.840 \pm 0.006$ & $70.8\%$ \\
\textbf{Officer's Arena (Hybrid)} & \textbf{BKT + HLR} & $\mathbf{0.862 \pm 0.004}$ & $\mathbf{89.6\%}$ \\
\bottomrule
\end{tabularx}
\end{table}

\subsection{Knowledge Tracing \& Memory Dynamics}
Table~\ref{tab:kt_benchmarks} demonstrates that the hybrid BKT-HLR framework achieves $\text{AUC-ROC} = 0.862$ and an $89.6\%$ detection accuracy for fragile misconception nodes, outperforming standard and deep KT architectures.

\begin{table}[h]
\centering
\caption{Inter-Rater Reliability for Mains Automated Essay Scoring.}
\label{tab:aes_agreement}
\small
\begin{tabularx}{\columnwidth}{lrrr}
\toprule
\textbf{Evaluator / Model Setup} & \textbf{QWK $\uparrow$} & \textbf{Pearson $r$ $\uparrow$} & \textbf{MAE $\downarrow$} \\
\midrule
\textbf{Human Rater 1 vs. Human Rater 2} & \textbf{0.871} & \textbf{0.884} & \textbf{4.12} \\
\textbf{Vanilla GPT-4o (Zero-Shot)} & $0.588$ & $0.612$ & $11.45$ \\
\textbf{Vanilla Claude 3.5 Sonnet} & $0.624$ & $0.648$ & $9.80$ \\
\textbf{LLaMA-3.3 70B + Basic RAG} & $0.738$ & $0.752$ & $7.15$ \\
\textbf{Officer's Arena AES (Gemini 2.5)} & $\mathbf{0.842}$ & $\mathbf{0.856}$ & $\mathbf{4.85}$ \\
\bottomrule
\end{tabularx}
\end{table}

\subsection{Mains Automated Essay Scoring (AES) Alignment}
As shown in Table~\ref{tab:aes_agreement}, unconstrained zero-shot models display low human alignment ($\text{QWK} = 0.588 - 0.624$). Officer's Arena's grounded 5D rubric pipeline achieves a Quadratic Weighted Kappa of $\text{QWK} = 0.842$ and Pearson $r = 0.856$, closely approaching the human benchmark ($\text{QWK} = 0.871$).

\begin{table}[h]
\centering
\caption{RAGAS Retrieval Quality Benchmarks over Canonical Vault.}
\label{tab:ragas_benchmarks}
\small
\begin{tabularx}{\columnwidth}{lrrr}
\toprule
\textbf{Metric} & \textbf{Target} & \textbf{Naive RAG} & \textbf{Officer's Arena Vault} \\
\midrule
\textbf{Context Precision} & $\ge 0.85$ & $0.684$ & $\mathbf{0.914}$ ($+23.0\%$) \\
\textbf{Context Recall} & $\ge 0.85$ & $0.712$ & $\mathbf{0.892}$ ($+18.0\%$) \\
\textbf{Faithfulness (Zero-Halluc.)} & $\ge 0.95$ & $0.785$ & $\mathbf{0.988}$ ($+20.3\%$) \\
\textbf{Answer Relevance} & $\ge 0.85$ & $0.742$ & $\mathbf{0.926}$ ($+18.4\%$) \\
\textbf{Mean Retrieval Latency} & $\le 50\text{ms}$ & $112\text{ms}$ & $\mathbf{14.2\text{ms}}$ ($7.9\times$ Speedup) \\
\bottomrule
\end{tabularx}
\end{table}

\begin{table*}[t]
\centering
\caption{Qualitative Case Study of Grounded Mains Automated Essay Evaluation (UPSC GS Paper II).}
\label{tab:case_study}
\small
\begin{tabularx}{\textwidth}{p{3.5cm}Xp{5.5cm}cc}
\toprule
\textbf{Evaluation Dimension} & \textbf{Candidate Submission Excerpt} & \textbf{Automated Feedback \& Grounded Citation} & \textbf{System} & \textbf{Human} \\
\midrule
\textbf{Directive Adherence ($s_1$)} & \multirow{5}{=}{``The Basic Structure doctrine was born in Kesavananda Bharati (1973)... Article 368 gives amending power but not power to destroy the Constitution. In Minerva Mills, SC struck down 42nd Amendment clauses. It balances power rather than creating supremacy.''} & Candidate balances both judicial review and legislative supremacy mandates without bias. & $8.5 / 10$ & $8.5 / 10$ \\
\textbf{Factual Grounding ($s_2$)} & & \textbf{Verbatim Citation}: \textit{``In Kesavananda Bharati (1973), the Supreme Court overruled Golaknath... Parliament cannot alter the basic structure.''} [Laxmikanth, Ch. 11, p. 11.2] & $9.0 / 10$ & $9.0 / 10$ \\
\textbf{PESTLE Breadth ($s_3$)} & & Covers constitutional governance, legal balance, and institutional checks and balances. & $8.0 / 10$ & $8.0 / 10$ \\
\textbf{Structural Flow ($s_4$)} & & Crisp tripartite structure with underlined landmark precedents. & $8.5 / 10$ & $8.0 / 10$ \\
\textbf{Way Forward ($s_5$)} & & Actionable synthesis on constitutionalism and Rule of Law. & $8.0 / 10$ & $8.5 / 10$ \\
\midrule
\textbf{TOTAL WEIGHTED} & \multicolumn{2}{l}{\textbf{Composite Scaled Score ($10 \times \sum w_d s_d$): High Precision Convergence ($\Delta = 0.5\%$)}} & \textbf{84.5} & \textbf{84.0} \\
\bottomrule
\end{tabularx}
\end{table*}

\begin{table}[h]
\centering
\caption{Infrastructure Throughput, Latency and Caching Performance.}
\label{tab:latency_benchmarks}
\small
\begin{tabularx}{\columnwidth}{lrrrr}
\toprule
\textbf{Microservice Operation} & \textbf{$p_{50}$ Latency} & \textbf{$p_{99}$ Latency} & \textbf{Cache Hit} & \textbf{Throughput} \\
\midrule
\textbf{3PL Next-Item Selection} & $\mathbf{4.2\text{ms}}$ & $\mathbf{12.1\text{ms}}$ & $99.4\%$ & $4,850\text{ req/s}$ \\
\textbf{EAP Quadrature Update} & $\mathbf{8.6\text{ms}}$ & $\mathbf{21.4\text{ms}}$ & Online & $2,420\text{ req/s}$ \\
\textbf{BKT State Transition} & $\mathbf{3.1\text{ms}}$ & $\mathbf{9.8\text{ms}}$ & $98.9\%$ & $5,600\text{ req/s}$ \\
\textbf{pgvector Dense Search} & $\mathbf{14.2\text{ms}}$ & $\mathbf{38.6\text{ms}}$ & $92.1\%$ & $1,150\text{ req/s}$ \\
\textbf{Mains AES (Gemini 2.5)} & $\mathbf{540\text{ms}}$ & $\mathbf{890\text{ms}}$ & Online & $125\text{ req/s}$ \\
\bottomrule
\end{tabularx}
\end{table}

\begin{table}[h]
\centering
\caption{System-Wide Component Ablation Matrix.}
\label{tab:ablation_matrix}
\small
\begin{tabularx}{\columnwidth}{lrrrr}
\toprule
\textbf{Configuration} & \textbf{$\Delta$ AUC} & \textbf{$\Delta$ Items} & \textbf{$\Delta$ QWK} & \textbf{$\Delta$ Faith.} \\
\midrule
\textbf{(A) Full System} & \textbf{0.864} & \textbf{57.8} & \textbf{0.842} & \textbf{98.8\%} \\
\textbf{(B) w/o EAP Quadrature} & $-0.045$ & $+10.7$ & --- & --- \\
\textbf{(C) w/o Max Fisher (MFI)} & $-0.098$ & $+28.6$ & --- & --- \\
\textbf{(D) w/o Half-Life Reg.} & $-0.032$ & --- & --- & --- \\
\textbf{(E) w/o Spatial VLM OCR} & $-0.024$ & --- & $-0.076$ & $-14.2\%$ \\
\textbf{(F) w/o 5D Rubric Grid} & --- & --- & $-0.218$ & $-12.6\%$ \\
\textbf{(G) w/o pgvector RAG} & --- & --- & $-0.254$ & $-20.3\%$ \\
\bottomrule
\end{tabularx}
\end{table}

\subsection{Component Ablation Studies}
As summarized in Table~\ref{tab:ablation_matrix}, removing EAP quadrature integration and MFI item selection increases candidate testing burden by over 39 items before achieving target measurement confidence. Removing the 5D rubric constraints causes the single largest decline in essay scoring agreement ($\Delta \text{QWK} = -0.218$), while naive OCR drops retrieval faithfulness by $14.2\%$.

% ==============================================================================
% SECTION 6: ETHICAL IMPLICATIONS, LIMITATIONS & FUTURE WORK
% ==============================================================================
\section{Ethical Implications, Limitations \& Future Work}

\subsection{Ethical Governance \& Fairness}
Officer's Arena implements continuous Differential Item Functioning (DIF) surveillance using Lord's $\chi^2$ Wald test:
\begin{equation}
    \chi_{\text{Lord}}^2 = (\hat{\boldsymbol{\beta}}_1 - \hat{\boldsymbol{\beta}}_2)^\top \left( \boldsymbol{\Sigma}_1 + \boldsymbol{\Sigma}_2 \right)^{-1} (\hat{\boldsymbol{\beta}}_1 - \hat{\boldsymbol{\beta}}_2)
    \label{eq:lords_chi2}
\end{equation}
Items exhibiting significant DIF ($\chi_{\text{Lord}}^2 > \chi_{\text{crit}}^2, p < 0.01$) are automatically quarantined to prevent systematic linguistic bias between Hindi and English cohorts. All interaction logs are pseudonymized via salted SHA-256 hashes with column-level database encryption \cite{Yan2024,Albadarin2024}.

\subsection{Limitations}
Automated evaluation of hand-drawn geographical sketch maps and flowchart diagrams in paper-written Mains answer scripts remains an open multimodal challenge. Furthermore, real-time numerical EAP quadrature at scale necessitates aggressive Redis caching of calibrated item vectors.

\subsection{Future Work: OmniGraph \& ChronoFact Engine}
Our Phase 4 roadmap introduces:
\begin{enumerate}[leftmargin=*]
    \item \textbf{OmniGraph Hypergraph}: A traversable semantic graph linking all 25,236 questions, distractor options, and 38 canonical textbooks across 114 KCs with prerequisite and contradiction edges \cite{Liang2024}.
    \item \textbf{ChronoFact Timeline Engine}: Multi-century causal event graphs structuring historical landmark sequences and policy evolutions.
    \item \textbf{Automated Item Generation (AIG)}: Synthesizing novel practice items with psychometric invariance matching target difficulty thresholds ($b_i = \theta_{\text{target}}$).
\end{enumerate}

% ==============================================================================
% SECTION 7: CONCLUSION
% ==============================================================================
\section{Conclusion}
In this paper, we introduced \textbf{Officer's Arena}, a comprehensive, mathematically grounded intelligent tutoring and psychometric assessment platform for high-stakes standardized examinations (UPSC CSE and CDS). By unifying 3PL-IRT Computerized Adaptive Testing ($42.2\%$ test length reduction), Bayesian Knowledge Tracing with Half-Life Regression memory decay ($89.6\%$ fragile misconception detection), vision-language multimodal document parsing (51.88M tokens indexed), and grounded 5D rubric essay scoring ($\text{QWK} = 0.842, 98.8\%$ citation faithfulness), Officer's Arena bridges the historical divide between classical psychometrics and generative foundation models for cognitive mastery.

% ==============================================================================
% REFERENCES
% ==============================================================================
\begin{thebibliography}{99}

\bibitem{Kim2022}
G.~Kim, T.~Hong, M.~Yim, J.~Nam, J.~Park, J.~Yim, W.~Hwang, S.~Yun, D.~Han, and S.~Park, ``OCR-free document understanding transformer,'' in \emph{Computer Vision -- ECCV 2022}, Springer, 2022, pp. 498--517. \url{https://doi.org/10.1007/978-3-031-20074-8_29}

\bibitem{Wang2024}
D.~Wang, N.~Raman, M.~Sibue, Z.~Ma, P.~Babkin, S.~Kaur, Y.~Pei, A.~Nourbakhsh, and X.~Liu, ``DocLLM: A layout-aware generative language model for multimodal document understanding,'' in \emph{Proc. 62nd Annu. Meeting Assoc. Comput. Linguistics (ACL 2024)}, 2024, pp. 8529--8548. \url{https://doi.org/10.18653/v1/2024.acl-long.463}

\bibitem{Zhang2024}
N.~Zhang, H.~Cheng, J.~Chen, Z.~Jiang, J.~Huang, Y.~Xue, and L.~Jin, ``M2Doc: A multi-modal fusion approach for document layout analysis,'' in \emph{Proc. AAAI Conf. Artif. Intell.}, vol.~38, no.~7, 2024, pp. 7233--7241. \url{https://doi.org/10.1609/aaai.v38i7.28552}

\bibitem{Abramovich2024}
O.~Abramovich, N.~Nayman, S.~Fogel, I.~Lavi, R.~Litman, S.~Tsiper, R.~Tichauer, S.~Appalaraju, S.~Mazor, and R.~Manmatha, ``VisFocus: Prompt-guided vision encoders for OCR-free dense document understanding,'' in \emph{Computer Vision -- ECCV 2024}, Springer, 2024, pp. 241--259. \url{https://doi.org/10.1007/978-3-031-73242-3_14}

\bibitem{Poznanski2025}
J.~Poznanski, J.~Borchardt, J.~Dunkelberger, R.~Huff, D.~Lin, A.~Rangapur, C.~Wilhelm, K.~Lo, and L.~Soldaini, ``olmOCR: Unlocking trillions of tokens in PDFs with vision language models,'' \emph{arXiv preprint arXiv:2502.18443}, 2025.

\bibitem{Zhuang2026}
Y.~Zhuang, Q.~Liu, H.~Bi, Z.~Huang, W.~Huang, J.~Li, J.~Yu, Z.~Liu, Z.~Hu, Y.~Hong, Z.~A.~Pardos, H.~Ma, M.~Zhu, S.~Wang, and E.~Chen, ``Survey of computerized adaptive testing: A machine learning perspective,'' \emph{IEEE Trans. Pattern Anal. Mach. Intell.}, vol.~48, no.~8, 2026, pp. 8744--8763. \url{https://doi.org/10.1109/TPAMI.2026.3672850}

\bibitem{Sharpnack2024}
J.~Sharpnack, P.~Mulcaire, K.~Bicknell, G.~LaFlair, and K.~Yancey, ``AutoIRT: Calibrating item response theory models with automated machine learning,'' \emph{arXiv preprint arXiv:2409.08823}, 2024.

\bibitem{Sharpnack2025}
J.~Sharpnack, K.~Hao, P.~Mulcaire, K.~Bicknell, G.~LaFlair, K.~Yancey, and A.~A.~von Davier, ``BanditCAT and AutoIRT: Machine learning approaches to computerized adaptive testing and item calibration,'' in \emph{Proc. Large Foundation Models for Educational Assessment (PMLR 264)}, 2025, pp. 121--135.

\bibitem{Shin2025}
H.~J.~Shin, C.~König, F.~Robin, A.~Frey, and K.~Yamamoto, ``Robustness of Item Response Theory models under the PISA multistage adaptive testing designs,'' \emph{J. Educ. Meas.}, vol.~62, no.~3, 2025, pp. 392--414. \url{https://doi.org/10.1111/jedm.12409}

\bibitem{Li2025}
J.~Li, J.~Sun, M.~Shao, Y.~Lai, and C.~Chen, ``A new multidimensional computerized testing approach: On-the-fly assembled multistage adaptive testing based on multidimensional item response theory,'' \emph{Mathematics}, vol.~13, no.~4, 2025, p. 594. \url{https://doi.org/10.3390/math13040594}

\bibitem{Ke2024}
F.~Ke, W.~Wang, W.~Tan, L.~Du, Y.~Jin, Y.~Huang, and H.~Yin, ``HiTSKT: A hierarchical transformer model for session-aware knowledge tracing,'' \emph{Knowl.-Based Syst.}, vol.~284, 2024, p. 111300. \url{https://doi.org/10.1016/j.knosys.2023.111300}

\bibitem{Liang2024}
Z.~Liang, R.~Wu, Z.~Liang, J.~Yang, L.~Wang, and J.~Su, ``GELT: A graph embeddings based lite-transformer for knowledge tracing,'' \emph{PLOS ONE}, vol.~19, no.~5, 2024, p. e0301714. \url{https://doi.org/10.1371/journal.pone.0301714}

\bibitem{Lu2024}
Y.~Lu, L.~Tong, and Y.~Cheng, ``Advanced knowledge tracing: Incorporating process data and curricula information via an attention-based framework for accuracy and interpretability,'' \emph{J. Educ. Data Mining}, vol.~16, no.~2, 2024, pp. 58--84. \url{https://doi.org/10.5281/zenodo.13712553}

\bibitem{BadrinathPardos2025}
A.~Badrinath and Z.~Pardos, ``Optimizing Bayesian Knowledge Tracing with neural network parameter generation,'' \emph{J. Educ. Data Mining}, vol.~17, no.~1, 2025, pp. 41--65. \url{https://doi.org/10.5281/zenodo.14707987}

\bibitem{Jung2025}
H.~Jung, J.~Yoo, Y.~Yoon, and Y.~Jang, ``CLST: Cold-start mitigation in knowledge tracing by aligning a generative language model as a students' knowledge tracer,'' \emph{J. Educ. Data Mining}, vol.~17, no.~2, 2025, pp. 86--117. \url{https://doi.org/10.5281/zenodo.17832627}

\bibitem{Yan2024}
L.~Yan, L.~Sha, L.~Zhao, Y.~Li, R.~Martinez-Maldonado, G.~Chen, X.~Li, Y.~Jin, and D.~Gašević, ``Practical and ethical challenges of large language models in education: A systematic scoping review,'' \emph{Brit. J. Educ. Technol.}, vol.~55, no.~1, 2024, pp. 90--112. \url{https://doi.org/10.1111/bjet.13370}

\bibitem{Albadarin2024}
Y.~Albadarin, M.~Saqr, N.~Pope, and M.~Tukiainen, ``A systematic literature review of empirical research on ChatGPT in education,'' \emph{Discover Educ.}, vol.~3, 2024, p. 60. \url{https://doi.org/10.1007/s44217-024-00138-2}

\bibitem{Hevia2025}
J.~S.~Hevia, F.~Arredondo, and V.~Kumar, ``Towards an efficient, customizable, and accessible AI tutor,'' in \emph{Proc. Innovation and Responsibility in AI-Supported Education Workshop (PMLR 273)}, 2025, pp. 250--254.

\bibitem{Kestin2025}
G.~Kestin, K.~Miller, A.~Klales, T.~Milbourne, and G.~Ponti, ``AI tutoring outperforms in-class active learning: An RCT introducing a novel research-based design in an authentic educational setting,'' \emph{Sci. Rep.}, vol.~15, 2025, p. 17458. \url{https://doi.org/10.1038/s41598-025-97652-6}

\bibitem{Bai2025}
H.~Bai, W.~C.~Lui, and P.~V.~Khiatani, ``Promoting student engagement with GPTutor: An intelligent tutoring system powered by generative AI,'' \emph{Int. J. Educ. Technol. High. Educ.}, vol.~22, 2025, p. 77. \url{https://doi.org/10.1186/s41239-025-00571-9}

\bibitem{Lord1980}
F.~M.~Lord, \emph{Applications of Item Response Theory to Practical Testing Problems}.\hskip 1em plus 0.5em minus 0.4em\relax Lawrence Erlbaum Associates, 1980.

\bibitem{Birnbaum1968}
A.~Birnbaum, ``Some latent trait models and their use in inferring an examinee's ability,'' in \emph{Statistical Theories of Mental Test Scores}, F.~M.~Lord and M.~R.~Novick, Eds.\hskip 1em plus 0.5em minus 0.4em\relax Addison-Wesley, 1968, pp. 397--479.

\bibitem{BockMislevy1982}
R.~D.~Bock and R.~J.~Mislevy, ``Adaptive EAP estimation of ability in a microcomputer environment,'' \emph{Appl. Psychol. Meas.}, vol.~6, no.~4, 1982, pp. 431--444. \url{https://doi.org/10.1177/014662168200600405}

\bibitem{BockAitkin1981}
R.~D.~Bock and M.~Aitkin, ``Marginal maximum likelihood estimation of item parameters: Application of an EM algorithm,'' \emph{Psychometrika}, vol.~46, no.~4, 1981, pp. 443--459. \url{https://doi.org/10.1007/BF02293801}

\bibitem{CorbettAnderson1995}
A.~T.~Corbett and J.~R.~Anderson, ``Knowledge tracing: Modeling the acquisition of procedural knowledge,'' \emph{User Model. User-Adapt. Interact.}, vol.~4, 1995, pp. 253--278. \url{https://doi.org/10.1007/BF01099821}

\bibitem{SettlesMeeder2016}
B.~Settles and B.~Meeder, ``A trainable spaced repetition model for language learning,'' in \emph{Proc. 54th Annu. Meeting Assoc. Comput. Linguistics (ACL 2016)}, 2016, pp. 1848--1858. \url{https://doi.org/10.18653/v1/P16-1174}

\bibitem{Wozniak1990}
P.~A.~Woźniak, ``Optimization of learning,'' Master's thesis, Univ. Technol., Poznań, Poland, 1990.

\bibitem{Lewis2020}
P.~Lewis, E.~Perez, A.~Piktus, F.~Petroni, V.~Karpukhin, N.~Goyal, H.~Küttler, M.~Lewis, W.~Yih, T.~Rocktäschel, S.~Riedel, and D.~Kiela, ``Retrieval-augmented generation for knowledge-intensive NLP tasks,'' in \emph{Adv. Neural Inf. Process. Syst. (NeurIPS)}, vol.~33, 2020, pp. 9459--9474.

\end{thebibliography}

\end{document}
